\documentclass{beamer}
\usepackage[T1]{fontenc}
\usepackage[utf8]{inputenc}
\usepackage[french]{babel}
\usepackage{lmodern}

\title{Resume LLM du dashboard energie}
\author{Projet Smart Building}
\date{}

\begin{document}

\begin{frame}{Comment fonctionne le resume LLM ?}
\small
\begin{columns}[T]
\begin{column}{0.49\textwidth}
\textbf{1. Ce que produit la chaine de calcul}
\begin{itemize}
\item Les previsions d'entree et l'etat courant sont envoyes a l'optimiseur.
\item L'optimiseur calcule un plan 24h :
\[
(P_{in}, P_{go}, PV, P_{ch}, P_{dis}, E_{bat}, S)
\]
\item Un module de reporting transforme ce plan en un \textit{payload} structure.
\end{itemize}

\vspace{0.4em}
\textbf{2. Ce que le LLM recupere vraiment}
\begin{itemize}
\item Totaux energie : demande servie, achat reseau, vente reseau, PV utilise, charge/decharge batterie.
\item Indicateurs derives : cout d'achat, revenu de vente, cout net, taux d'autoconsommation PV, taux de service de la flexibilite.
\item Infos temporelles : heures et amplitudes des pics d'import, d'export, de charge et de decharge.
\item Comparaison \textit{baseline} vs optimiseur : cout et emissions CO$_2$.
\end{itemize}
\end{column}

\begin{column}{0.49\textwidth}
\textbf{3. Ce qu'on lui donne sur le modele d'optimisation}
\begin{itemize}
\item Horizon de decision : 24 heures.
\item Fonction objectif : minimiser le cout energie total en tenant compte de l'achat reseau, de la vente, du PV, de la batterie et de la charge flexible.
\item Variables de decision : import/export reseau, charge/decharge batterie, energie stockee, service flexible.
\item Contraintes : bilan de puissance, dynamique de batterie, bornes de puissance, capacite, non-simultaneite charge/decharge.
\item Parametres : limites reseau, capacite batterie, penalites et coefficients economiques/carbone.
\end{itemize}

\vspace{0.4em}
\textbf{4. Role du LLM}
\begin{itemize}
\item Il ne voit pas directement les graphes ni les donnees brutes.
\item Il recoit uniquement ce resume chiffre + la description du modele.
\item Il redige ensuite une explication naturelle : \og pourquoi l'optimiseur a choisi ce plan\fg{} au regard de l'objectif et des contraintes.
\item Si l'API LLM est indisponible, un resume deterministe de secours est genere.
\end{itemize}
\end{column}
\end{columns}
\end{frame}

\end{document}
